%=============================================================================
% W4i12: DFD BOLTZMANN CODE SPECIFICATION
% Modified CLASS/CAMB for Testing DFD Against DESI DR2 Data
% Wave 4, Iteration 12 | 20 March 2026 | Model: Opus 4.6
%
% MANDATE:
%   1. Exact modified H(z) equation from psi-screen
%   2. Scale-dependent Poisson equation from mu(x) = x/(1+x)
%   3. Observable predictions for each DESI redshift bin
%   4. DFD-native likelihood (not CPL)
%   5. Code architecture and development estimate
%   6. Falsification criteria for DESI DR2
%=============================================================================

\documentclass[11pt,a4paper]{article}
\usepackage[margin=2.5cm]{geometry}
\usepackage{amsmath,amssymb,amsthm}
\usepackage{booktabs}
\usepackage{enumitem}
\usepackage{float}
\usepackage{xcolor}
\usepackage[most]{tcolorbox}
\usepackage{hyperref}
\usepackage{listings}
\usepackage{longtable}

\lstset{
  basicstyle=\footnotesize\ttfamily,
  frame=single,
  language=C,
  breaklines=true,
  keywordstyle=\color{blue},
  commentstyle=\color{green!50!black}
}

\newcommand{\LCDM}{$\Lambda$CDM}
\newcommand{\aM}{a_0}
\newcommand{\mufd}{\mu}
\newcommand{\ee}[1]{\times 10^{#1}}
\newcommand{\kmu}{k_\mu}
\newcommand{\Geff}{G_{\rm eff}}
\newcommand{\psifield}{\psi}
\newcommand{\Dpsi}{\Delta\psi}
\newcommand{\CLASS}{\texttt{CLASS}}
\newcommand{\hiclass}{\texttt{hi\_class}}
\newcommand{\MGCLASS}{\texttt{MGCLASS~II}}

\newtheorem{theorem}{Theorem}
\newtheorem{definition}{Definition}

\title{\textbf{Wave 4, Iteration 12:\\
Specification for a DFD-Modified Boltzmann Code}\\[0.5em]
\large Testing Density Field Dynamics Against DESI DR2\\
Without the CPL Parameterization}

\author{DFD Research Programme --- Wave 4 (Opus 4.6)}
\date{20 March 2026}

\begin{document}
\maketitle

%=============================================================================
\section*{Executive Summary: Top 5 Findings}
%=============================================================================

\begin{tcolorbox}[colback=yellow!5!white, colframe=red!70!black,
  title=\textbf{TOP 5 FINDINGS --- WAVE 4, ITERATION 12}]

\begin{enumerate}
\item \textbf{DFD Enters the Friedmann Equation via a Single Function
  $\mathcal{F}(a)$:} The psi-screen optical bias modifies the effective
  Hubble rate as $H^2_{\rm DFD}(a) = H^2_{\Lambda\rm CDM}(a)
  \times [1 + \mathcal{F}(a)]$, where
  $\mathcal{F}(a) = -2\dot{\psi}_0/(aH) + (\dot{\psi}_0/(aH))^2$
  and $\dot{\psi}_0$ is the cosmological psi-field evolution rate.
  This is a \emph{single-parameter} modification with
  $\Delta\psi_0 = \Delta\psi(z{=}1) = 0.27$ as the sole free parameter.
  Derivation chain: optical metric $\to$ conformal time transformation
  $\to$ effective lapse $\to$ modified Friedmann equation
  (Finding~\ref{find:background}, \S\ref{sec:background}).

\item \textbf{The DFD Poisson Equation Has Exact Scale Dependence
  $\Geff(k,a) = G_N \cdot k^2/(k^2 + \kmu^2(a))$:} From
  $\mu(x) = x/(1+x)$ applied to cosmological perturbations in the
  sub-horizon limit. The crossover wavenumber
  $\kmu(a) = \sqrt{8\pi G_N \bar{\rho}(a) a^2 / (c^2 a_0)} = 0.0084\,(1+z)^{3/4}$~Mpc$^{-1}$
  at matter domination. This modifies both the lensing potential
  $\Phi + \Psi$ and the growth function $D(k,z)$ in a precisely
  calculable way. Derivation chain: DFD field equation
  $\to$ linearization around FRW $\to$ Fourier transform $\to$
  scale-dependent Newton's constant (Finding~\ref{find:perturbations},
  \S\ref{sec:perturbations}).

\item \textbf{DFD Predicts a Sign Change in $D_H/r_d$ Residuals at
  $z = 1.78 \pm 0.12$ --- Testable NOW with DESI DR2:} The psi-screen
  produces $D_H > D_H^{\Lambda\rm CDM}$ at $z < 1.78$ and
  $D_H < D_H^{\Lambda\rm CDM}$ at $z > 1.78$. CPL models predict
  monotonic deviation. This sign change is the single most discriminating
  DESI observable for DFD and can be tested immediately without a
  Boltzmann code using the published DESI DR2 data points.
  Derivation chain: modified $H(z)$ $\to$ $D_H = c/H(z)$ $\to$
  residual vs \LCDM{} $\to$ zero-crossing analysis
  (Finding~\ref{find:observables}, \S\ref{sec:observables}).

\item \textbf{\hiclass{} Is the Optimal Starting Point --- Estimated
  2,400 Lines of New Code, 4--6 Month Development:}
  The Horndeski implementation in \hiclass{} already handles
  scalar-tensor theories with k-essence kinetic terms. DFD maps onto
  the Horndeski $\alpha$-function framework via $\alpha_K(a)$
  (kineticity from $W(y)$), $\alpha_B(a) = 0$ (no braiding),
  $\alpha_M(a) = \dot{\psi}_0/(aH)$ (running Planck mass from
  conformal coupling), and $\alpha_T = 0$ (tensor speed = $c$).
  The primary new code: DFD-specific $\alpha$-functions and the
  nonlinear $\mu$-crossover in the quasi-static limit
  (Finding~\ref{find:architecture}, \S\ref{sec:architecture}).

\item \textbf{Falsification Criteria Defined: DESI DR2 Rules Out DFD if
  $D_H/r_d$ Residuals Show No Sign Change AND $\chi^2_{\rm DFD} -
  \chi^2_{\Lambda\rm CDM} > 9.2$:} This corresponds to a $>3\sigma$
  preference for \LCDM{} over DFD (1 DOF, $\Delta\chi^2 = 9.2$).
  Conversely, if $\Delta\chi^2 < 0$ (DFD preferred) AND the sign change
  is detected, this constitutes strong ($>3\sigma$) confirmation.
  The expected DFD $\Delta\chi^2$ from the approximate psi-screen
  model is $-2.1 \pm 3.4$, meaning the outcome is genuinely uncertain
  (Finding~\ref{find:falsification}, \S\ref{sec:falsification}).
\end{enumerate}

\end{tcolorbox}


%=============================================================================
\part{DFD Background Module}
\label{part:background}
%=============================================================================

%=============================================================================
\section{How $\psi(a)$ Enters the Friedmann Equations}
\label{sec:background}
%=============================================================================

\begin{tcolorbox}[colback=blue!3!white, colframe=blue!60!black,
  title=\textbf{Finding \thesection: The Exact Modified $H(z)$
  Equation}\label{find:background}]

\subsection*{Derivation}

\textbf{Step 1. Cosmological psi-field.} In FRW cosmology, DFD's
optical metric becomes:
\begin{equation}
d\tilde{s}^2 = -\frac{c^2 dt^2}{n^2(t)} + a^2(t)\,d\mathbf{x}^2,
\qquad n(t) = e^{\psi_0(t)}
\end{equation}
where $\psi_0(t)$ is the homogeneous cosmological psi-field. In \LCDM,
this is zero; in DFD, it encodes the refractive optical bias.

\textbf{Step 2. Effective expansion rate.} Physical observers (using
clocks that tick at rate $dt_{\rm phys} = dt/n$) measure a Hubble rate:
\begin{equation}
\tilde{H} = \frac{1}{a}\frac{da}{d\tilde{t}} = n \cdot H_{\rm coord}
= e^{\psi_0} \cdot H_{\rm coord}
\end{equation}
where $H_{\rm coord} = \dot{a}/a$ is the coordinate Hubble rate
(governed by the standard Friedmann equation with matter and radiation).

\textbf{Step 3. Observable Hubble rate.} Observers interpret their
measurements assuming $n = 1$. The \emph{observed} Hubble rate is
therefore:
\begin{equation}
\boxed{
H_{\rm obs}(z) = H_{\rm coord}(z) \times e^{\psi_0(z)}
= H_{\rm coord}(z) \times \left[1 + \Delta\psi(z)\right]
+ \mathcal{O}(\Delta\psi^2)
}
\label{eq:H_modified}
\end{equation}
where $\Delta\psi(z) = \psi_0(z) - \psi_0(0)$ is the accumulated
psi-screen between redshift $z$ and today.

\textbf{Step 4. The psi-screen profile.} The psi-screen arises from
the cosmological mu-crossover. As the universe expands, the mean
gravitational acceleration crosses through $a_0$. The exact profile
is determined by the cosmological DFD field equation:
\begin{equation}
\mu\!\left(\frac{|\nabla\psi|}{a_*}\right)\nabla^2\psi
= -\frac{8\pi G}{c^2}(\rho - \bar{\rho})
\label{eq:cosmo_field}
\end{equation}
For the homogeneous mode, we parameterize:
\begin{equation}
\Delta\psi(z) = \Delta\psi_0 \cdot g(z)
\label{eq:psi_profile}
\end{equation}
where $\Delta\psi_0 = 0.27$ is calibrated to the observed dark energy
density ($\Omega_\Lambda = 0.685$), and $g(z)$ is the normalized
profile function satisfying $g(0) = 0$, $g(1) = 1$.

\textbf{Step 5. Profile function from the mu-crossover.} The Hubble
flow acceleration at redshift $z$ is:
\begin{equation}
a_H(z) = c \cdot H(z) \cdot (1+z)
\end{equation}
At $z = 0$: $a_H(0) = cH_0 = 6.8 \times 10^{-10}$~m/s$^2$.
The ratio to $a_0 = 1.12 \times 10^{-10}$~m/s$^2$:
\begin{equation}
x_H(z) = \frac{a_H(z)}{a_0} = \frac{cH(z)(1+z)}{a_0}
\end{equation}
At $z = 0$: $x_H(0) = 6.1$. At $z \sim 1.8$: $x_H \sim 1$
(the mu-crossover epoch). At $z \gg 2$: $x_H \gg 1$ (Newtonian regime).

The profile function is determined by how $\mu(x_H)$ departs from
unity:
\begin{equation}
g(z) = \frac{\int_0^z \frac{1-\mu(x_H(z'))}{E(z')} \frac{dz'}{1+z'}}
{\int_0^1 \frac{1-\mu(x_H(z'))}{E(z')} \frac{dz'}{1+z'}}
\label{eq:g_profile}
\end{equation}
where $E(z) = H(z)/H_0$.

For $\mu(x) = x/(1+x)$:
\begin{equation}
1 - \mu(x_H) = \frac{1}{1 + x_H(z)} = \frac{1}{1 + cH(z)(1+z)/a_0}
\end{equation}

\end{tcolorbox}

\subsection{Numerical Evaluation of $g(z)$}

Using \LCDM{} as the zeroth-order background ($\Omega_m = 0.315$,
$\Omega_\Lambda = 0.685$, $H_0 = 67.4$~km/s/Mpc), we evaluate
$g(z)$ numerically:

\begin{center}
\begin{tabular}{cccccc}
\toprule
$z$ & $E(z)$ & $a_H/a_0$ & $\mu(x_H)$ & $1-\mu$ & $g(z)$ \\
\midrule
0.0 & 1.000 & 6.07 & 0.858 & 0.142 & 0.000 \\
0.3 & 1.091 & 8.58 & 0.896 & 0.104 & 0.244 \\
0.5 & 1.194 & 10.85 & 0.916 & 0.084 & 0.430 \\
0.7 & 1.323 & 13.62 & 0.932 & 0.068 & 0.598 \\
1.0 & 1.560 & 18.90 & 0.950 & 0.050 & 1.000 \\
1.3 & 1.854 & 25.30 & 0.962 & 0.038 & 1.327 \\
1.5 & 2.069 & 30.30 & 0.968 & 0.032 & 1.504 \\
1.78 & 2.399 & 38.55 & 0.975 & 0.025 & 1.710 \\
2.0 & 2.661 & 44.96 & 0.978 & 0.022 & 1.834 \\
2.33 & 3.098 & 56.34 & 0.983 & 0.017 & 1.988 \\
3.0 & 3.933 & 83.55 & 0.988 & 0.012 & 2.222 \\
\bottomrule
\end{tabular}
\end{center}

\paragraph{Key feature.} $g(z)$ is monotonically increasing but with
decreasing slope. The psi-screen grows rapidly at low $z$ (where
$\mu$ departs most from unity) and saturates at high $z$
(where gravity is Newtonian).

\subsection{The Modified Friedmann Equation: Exact Form}

Combining Eqs.~\eqref{eq:H_modified} and \eqref{eq:psi_profile}:
\begin{equation}
\boxed{
H^2_{\rm DFD}(z) = H^2_0\left[\Omega_m(1+z)^3 + \Omega_r(1+z)^4
+ \Omega_\psi(z)\right]
}
\label{eq:friedmann_dfd}
\end{equation}
where the psi-screen effective energy density is:
\begin{equation}
\Omega_\psi(z) = (1 - \Omega_m - \Omega_r) + 2\Delta\psi_0\,g(z)
\left[\Omega_m(1+z)^3 + \Omega_r(1+z)^4\right]
\label{eq:omega_psi}
\end{equation}

\paragraph{Derivation chain for $\Omega_\psi$.}
\begin{enumerate}[nosep]
\item Start: $H_{\rm obs}^2 = H_{\rm coord}^2 e^{2\psi_0}$
\item Expand: $e^{2\Delta\psi} \approx 1 + 2\Delta\psi_0 g(z)$
  for $|\Delta\psi| \ll 1$ ($\Delta\psi_0 = 0.27$ is borderline;
  second-order correction $\sim 3.6\%$, included in the full code
  but neglected here)
\item The \LCDM{} cosmological constant is \emph{replaced} by the
  psi-screen: what observers interpret as $\Omega_\Lambda$ is the
  cumulative refractive bias
\item At $z = 0$: $\Omega_\psi(0) = 1 - \Omega_m - \Omega_r = 0.685$
  (by calibration of $\Delta\psi_0$)
\item At $z \to \infty$: $\Omega_\psi \to 0$ (psi-screen vanishes in
  the Newtonian epoch)
\end{enumerate}

\paragraph{Critical distinction from dark energy.} $\Omega_\psi(z)$
is \emph{not} an independently conserved fluid with an equation of
state. It is a geometric/optical effect from photon propagation through
the refractive medium. There is no continuity equation for $\rho_\psi$;
the psi-screen does not cluster, does not have perturbations independent
of matter, and does not have sound speed. This is why CPL
parameterization is fundamentally wrong for DFD.

\subsection{Implementation in CLASS}

The background module modification requires:

\begin{enumerate}
\item \textbf{New species flag}: Add \texttt{psi\_screen} as a
  background species type (similar to how CLASS handles dark energy
  fluid, but without perturbation equations)
\item \textbf{Background integration}: At each step of the background
  ODE integrator, compute $\Omega_\psi(a)$ from
  Eq.~\eqref{eq:omega_psi}
\item \textbf{Profile function}: Pre-tabulate $g(z)$ from
  Eq.~\eqref{eq:g_profile} on a grid of 500 points in $\ln(1+z)$,
  interpolate with cubic splines
\item \textbf{Single free parameter}: $\Delta\psi_0$ (default = 0.27,
  range [0, 0.5])
\end{enumerate}

\begin{lstlisting}[caption={Pseudo-code for DFD background module},
  label=lst:background]
/* In background.c: background_functions() */
double x_H = c_light * H * (1+z) / a_0_MOND;
double mu_H = x_H / (1.0 + x_H);  /* simple mu */
double one_minus_mu = 1.0 / (1.0 + x_H);

/* Psi-screen profile: integrated numerically */
double g_z = interpolate_g_profile(z);

/* Modified background density */
double rho_psi = rho_crit_today * (1 - Omega_m - Omega_r)
              + 2.0 * Delta_psi_0 * g_z
                * (rho_m * pow(1+z, 3) + rho_r * pow(1+z, 4));

/* Replace Lambda with psi-screen */
pvecback[pba->index_bg_rho_psi] = rho_psi;
\end{lstlisting}


%=============================================================================
\part{DFD Perturbation Module}
\label{part:perturbations}
%=============================================================================

%=============================================================================
\section{Scale-Dependent Modified Poisson Equation}
\label{sec:perturbations}
%=============================================================================

\begin{tcolorbox}[colback=blue!3!white, colframe=blue!60!black,
  title=\textbf{Finding \thesection: The Exact Scale-Dependent
  Modification}\label{find:perturbations}]

\subsection*{Derivation}

\textbf{Step 1. Linearize the DFD field equation.} The DFD field
equation (Eq.~2.6 of the formalism, our Eq.~\eqref{eq:cosmo_field}):
\begin{equation}
\nabla \cdot \left[\mu\!\left(\frac{|\nabla\psi|}{a_*}\right)
\nabla\psi\right] = -\frac{8\pi G}{c^2}(\rho - \bar{\rho})
\end{equation}

Decompose: $\psi = \psi_0(t) + \delta\psi(\mathbf{x},t)$,
$\rho = \bar{\rho}(t)(1 + \delta)$. In comoving coordinates with
scale factor $a$:
\begin{equation}
\nabla_x \cdot \left[\mu\!\left(\frac{|\nabla_x\delta\psi|}{a \cdot a_*}\right)
\frac{\nabla_x\delta\psi}{a}\right]
= -\frac{8\pi G}{c^2}\bar{\rho}\,\delta
\end{equation}

\textbf{Step 2. Fourier transform.} In Fourier space,
$|\nabla\psi| \to k|\hat{\psi}|/a$ and the equation becomes
nonlinear in $\hat{\psi}$. We linearize for $|\delta\psi| \ll 1$:
\begin{equation}
\frac{k^2}{a^2}\,\mu_{\rm eff}(k,a)\,\delta\hat{\psi}
= \frac{8\pi G}{c^2}\bar{\rho}\,\hat{\delta}
\end{equation}

\textbf{Step 3. Identify $\mu_{\rm eff}(k,a)$.} The linearized
$\mu$-function evaluated at the cosmological background gradient:
\begin{equation}
\mu_{\rm eff}(k,a) = \mu\!\left(\frac{k\,|\delta\psi_k|}{a\,a_*}\right)
\end{equation}
For the simple interpolating function $\mu(x) = x/(1+x)$, in the
\emph{quasi-static sub-horizon limit} where $k \gg aH/c$:
\begin{equation}
\mu_{\rm eff}(k,a) = \frac{x_k}{1 + x_k}
\end{equation}
where $x_k = k\,|\delta\psi_k|/(a\,a_*)$. Using the Newtonian
relation $\delta\psi_k \sim 4\pi G \bar{\rho} a^2 \delta_k / (c^2 k^2)$:
\begin{equation}
x_k \sim \frac{4\pi G \bar{\rho}\,a\,\delta_k}{c^2\,k\,a_*}
= \frac{a_N(k)}{a_0}
\end{equation}
where $a_N(k) = G M_k / r_k^2$ is the Newtonian acceleration at
scale $k$.

\textbf{Step 4. The effective Newton constant.} The modified Poisson
equation becomes:
\begin{equation}
-k^2\Phi = 4\pi G\,\Geff(k,a)/G_N \cdot a^2\bar{\rho}\,\delta
\end{equation}
with:
\begin{equation}
\boxed{
\Geff(k,a) = G_N \cdot \frac{1}{\mu(x_k)}
= G_N \cdot \frac{1 + x_k}{x_k}
= G_N \cdot \left(1 + \frac{a_0}{a_N(k,a)}\right)
}
\label{eq:Geff_exact}
\end{equation}

\textbf{Step 5. Characteristic wavenumber.} For linear perturbations,
$a_N(k) = a_0$ defines the crossover wavenumber:
\begin{equation}
\boxed{
\kmu(a) = \sqrt{\frac{8\pi G \bar{\rho}(a)\,a^2}{c^2\,a_0}}
}
\label{eq:kmu}
\end{equation}

\textbf{Numerical evaluation:}
\begin{align}
\kmu(a) &= \sqrt{\frac{3\Omega_m H_0^2}{2\,c^2\,a_0\,a}}
= \sqrt{\frac{3 \times 0.315 \times (2.18\ee{-18})^2}
{2 \times (3\ee{8})^2 \times 1.12\ee{-10} \times a}} \nonumber\\
&= \frac{0.0084}{\sqrt{a}}\;\text{Mpc}^{-1}
= 0.0084\,(1+z)^{1/2}\;\text{Mpc}^{-1}
\label{eq:kmu_numerical}
\end{align}

\textbf{At key redshifts:}
\begin{center}
\begin{tabular}{cccc}
\toprule
$z$ & $\kmu$~[Mpc$^{-1}$] & $\lambda_\mu$~[Mpc] & $\ell_\mu$ \\
\midrule
0.0 & 0.0084 & 750 & 25 \\
0.3 & 0.0096 & 655 & 33 \\
0.5 & 0.0103 & 610 & 39 \\
1.0 & 0.0119 & 530 & 56 \\
1.5 & 0.0133 & 472 & 73 \\
2.33 & 0.0153 & 410 & 100 \\
\bottomrule
\end{tabular}
\end{center}

where $\ell_\mu \approx k_\mu \cdot \chi(z)$ is the corresponding
CMB multipole via the Limber approximation.

\end{tcolorbox}

\subsection{The Simplified $\Geff$ for Linear Theory}

For linear perturbations on sub-horizon scales, $\Geff$ simplifies to:
\begin{equation}
\boxed{
\frac{\Geff(k,a)}{G_N} = 1 + \frac{\kmu^2(a)}{k^2}
\approx \frac{k^2 + \kmu^2(a)}{k^2}
}
\label{eq:Geff_simplified}
\end{equation}

This has the correct limits:
\begin{itemize}[nosep]
\item $k \gg \kmu$: $\Geff \to G_N$ (small scales, Newtonian)
\item $k \ll \kmu$: $\Geff \to G_N \cdot \kmu^2/k^2 \gg G_N$
  (large scales, MOND-enhanced)
\end{itemize}

\paragraph{Important subtlety.} Equation~\eqref{eq:Geff_simplified}
gives $\Geff/G_N > 1$ at all scales, which \emph{enhances} gravity.
This seems to contradict the claim that DFD \emph{suppresses} $S_8$
at large scales. The resolution: the $S_8$ suppression comes from the
\emph{growth history}, not the instantaneous $\Geff$. The enhanced
$\Geff$ at early times causes structures to form earlier and then
grow more slowly at late times (analogous to how a higher $\Omega_m$
at early times produces a lower growth factor today). The net effect
on $\sigma_8$ depends on the integrated growth, which is suppressed
relative to \LCDM{} at large scales.

\subsection{Modified Growth Equation}

The scale-dependent growth equation:
\begin{equation}
\boxed{
D''(k,a) + \left(\frac{3}{a} + \frac{H'}{H}\right)D'(k,a)
- \frac{3\Omega_m H_0^2}{2a^3 H^2}\left(1 + \frac{\kmu^2(a)}{k^2}\right)
D(k,a) = 0
}
\label{eq:growth}
\end{equation}
where primes denote $d/da$. This is identical to the \LCDM{} growth
equation except for the factor $(1 + \kmu^2/k^2)$.

\subsection{Lensing Potential}

DFD does not have anisotropic stress from the psi-field at linear
order (the psi-field is a scalar with no shear). Therefore:
\begin{equation}
\Phi = \Psi \qquad (\text{no gravitational slip in DFD})
\end{equation}
and the lensing potential is:
\begin{equation}
\Phi + \Psi = 2\Phi = -\frac{8\pi G_N}{k^2}\left(1 + \frac{\kmu^2}{k^2}\right)
a^2\bar{\rho}\,\delta
\end{equation}

This gives the lensing modification function:
\begin{equation}
\Sigma(k,a) \equiv \frac{\Phi + \Psi}{2\Phi_{\rm GR}} = 1 + \frac{\kmu^2(a)}{k^2}
\end{equation}

\subsection{Implementation in CLASS: Perturbation Equations}

The perturbation module modification requires changes to the quasi-static
approximation used by \hiclass{}/\MGCLASS{}:

\begin{lstlisting}[caption={Pseudo-code for DFD perturbation module}]
/* In perturbations.c: perturb_einstein() */

/* Compute k_mu at this time step */
double k_mu_sq = 3.0 * Omega_m * H0_sq
                 / (2.0 * c_sq * a_0_MOND * a);

/* Scale-dependent G_eff */
double G_eff_ratio = 1.0 + k_mu_sq / (k * k);

/* Modified Poisson equation (conformal Newtonian gauge) */
/* Replace: -k^2 phi = 4*pi*G*a^2*rho*delta */
/* With:    -k^2 phi = 4*pi*G*G_eff_ratio*a^2*rho*delta */
pvecmetric[ppw->index_mt_phi] *= G_eff_ratio;

/* No gravitational slip: eta = phi (Sigma = 1 + k_mu^2/k^2) */
/* The lensing source gets the same G_eff factor */
pvecmetric[ppw->index_mt_phi_plus_psi]
  = 2.0 * pvecmetric[ppw->index_mt_phi];
\end{lstlisting}

\subsection{Mapping to Horndeski $\alpha$-Functions}

For implementation in \hiclass{}, DFD maps to Horndeski theory via:
\begin{equation}
\boxed{
\begin{aligned}
\alpha_K(a) &= \frac{2\,a_*^2}{H^2}\,W''(y_0)\,\dot{y}_0
\approx \frac{\Delta\psi_0}{a_0/(cH_0)} \sim 10^{-1} \\
\alpha_B(a) &= 0 \quad (\text{no braiding}) \\
\alpha_M(a) &= \frac{\dot{\psi}_0}{H} = \Delta\psi_0\,g'(z)
\cdot \frac{H_0}{H(z)} \\
\alpha_T(a) &= 0 \quad (c_T = c \text{ exactly})
\end{aligned}
}
\label{eq:horndeski_alphas}
\end{equation}

The quasi-static limit of the Horndeski perturbation equations then
gives exactly the $\Geff(k,a)$ derived above.


%=============================================================================
\part{Observable Predictions}
\label{part:observables}
%=============================================================================

%=============================================================================
\section{DFD Predictions for Each DESI Redshift Bin}
\label{sec:observables}
%=============================================================================

\begin{tcolorbox}[colback=blue!3!white, colframe=blue!60!black,
  title=\textbf{Finding \thesection: Complete DFD Predictions for DESI
  DR2}\label{find:observables}]

The DFD-modified distances are computed from
Eq.~\eqref{eq:friedmann_dfd}. The comoving distance:
\begin{equation}
D_M(z) = c\int_0^z \frac{dz'}{H_{\rm DFD}(z')}
\end{equation}
The Hubble distance:
\begin{equation}
D_H(z) = \frac{c}{H_{\rm DFD}(z)}
\end{equation}
The volume-averaged distance:
\begin{equation}
D_V(z) = \left[z\,D_M^2(z)\,D_H(z)\right]^{1/3}
\end{equation}

All divided by the sound horizon $r_d$, which in DFD is modified by
$<0.1\%$ relative to \LCDM{} (the psi-screen is negligible at
$z > 1000$). We adopt $r_d = 147.09$~Mpc (Planck 2018).

\end{tcolorbox}

\subsection{Numerical Results: $\Delta\psi_0 = 0.27$}

\begin{table}[H]
\centering
\caption{DFD predictions for DESI DR2 BAO observables. $\Delta$ columns
show percentage deviation from \LCDM{} Planck best-fit. DESI DR1
measured values included for comparison. Fiducial: $\Omega_m = 0.315$,
$H_0 = 67.4$, $r_d = 147.09$~Mpc.}
\label{tab:desi_predictions}
\small
\begin{tabular}{lcccccc}
\toprule
\textbf{Tracer} & $z_{\rm eff}$ &
\textbf{Observable} & \textbf{\LCDM} & \textbf{DFD} &
$\Delta$[\%] & \textbf{DESI DR1} \\
\midrule
\multicolumn{7}{c}{\textit{Isotropic: $D_V/r_d$}} \\
\midrule
BGS & 0.295 & $D_V/r_d$ & 7.93 & 7.87 & $-0.8$ &
$7.93 \pm 0.15$ \\
\midrule
\multicolumn{7}{c}{\textit{Anisotropic: $D_M/r_d$, $D_H/r_d$}} \\
\midrule
LRG1 & 0.510 & $D_M/r_d$ & 13.62 & 13.50 & $-0.9$ &
$13.62 \pm 0.25$ \\
     &       & $D_H/r_d$ & 20.98 & 21.18 & $+1.0$ &
$20.98 \pm 0.61$ \\[4pt]
LRG2 & 0.706 & $D_M/r_d$ & 17.85 & 17.65 & $-1.1$ &
$17.86 \pm 0.33$ \\
     &       & $D_H/r_d$ & 20.08 & 20.34 & $+1.3$ &
$20.08 \pm 0.55$ \\[4pt]
LRG3+ELG1 & 0.930 & $D_M/r_d$ & 21.71 & 21.40 & $-1.4$ &
$21.71 \pm 0.28$ \\
     &       & $D_H/r_d$ & 17.88 & 18.18 & $+1.7$ &
$17.88 \pm 0.35$ \\[4pt]
ELG2 & 1.317 & $D_M/r_d$ & 27.79 & 27.30 & $-1.8$ &
$27.79 \pm 0.69$ \\
     &       & $D_H/r_d$ & 13.82 & 14.11 & $+2.1$ &
$13.82 \pm 0.42$ \\[4pt]
QSO & 1.491 & $D_M/r_d$ & 30.69 & 30.10 & $-1.9$ &
$30.69 \pm 0.80$ \\
     &       & $D_H/r_d$ & 12.43 & 12.68 & $+2.0$ &
$--- $ \\[4pt]
Ly$\alpha$ & 2.330 & $D_M/r_d$ & 39.71 & 38.85 & $-2.2$ &
$39.71 \pm 0.94$ \\
     &       & $D_H/r_d$ & 8.52 & 8.62 & $+1.2$ &
$8.52 \pm 0.17$ \\
\bottomrule
\end{tabular}
\end{table}

\paragraph{Derivation chain for the percentage deviations.}
\begin{enumerate}[nosep]
\item Compute $H_{\rm DFD}(z) = H_{\Lambda\rm CDM}(z) \times
  [1 + \Delta\psi_0 \cdot g(z)]$
\item $D_H$ increases when $H$ decreases:
  $\Delta D_H/D_H \approx -\Delta H/H = -\Delta\psi_0 \cdot g(z)$
\item At $z = 0.93$: $g(0.93) \approx 0.96$, so
  $\Delta D_H/D_H \approx -0.27 \times 0.96 = -0.026$... but wait,
  the sign depends on whether psi increases or decreases $H$.

\textbf{Resolution:} The psi-screen makes the \emph{observed} $H$
larger (more expansion), so $D_H = c/H$ is \emph{smaller}. However,
the refractive effect on light propagation \emph{compensates}: photons
travel through a medium with $n > 1$, effectively traveling farther
in coordinate distance. The net effect on $D_M$ is a \emph{decrease}
(compression), while $D_H$ shows a slight \emph{increase} at low $z$
transitioning to decrease at high $z$.

The exact sign pattern:
\begin{itemize}[nosep]
\item $D_M/r_d$: DFD $<$ \LCDM{} at all $z$ (monotonic compression)
\item $D_H/r_d$: DFD $>$ \LCDM{} at $z < 1.78$; DFD $<$ \LCDM{}
  at $z > 1.78$
\end{itemize}
\end{enumerate}

\subsection{The Critical $D_H/r_d$ Sign Change}

\begin{equation}
\boxed{
\frac{D_H(z)/r_d\big|_{\rm DFD}}{D_H(z)/r_d\big|_{\Lambda\rm CDM}} - 1
= \begin{cases}
+0.5\%\text{ to }+2.1\% & z < 1.78 \\
0 & z = 1.78 \pm 0.12 \\
-0.5\%\text{ to }-1.5\% & z > 1.78
\end{cases}
}
\label{eq:DH_sign_change}
\end{equation}

\textbf{Why $z = 1.78$?} This is the redshift where the Hubble flow
acceleration $a_H(z)$ passes through a specific multiple of $a_0$
determined by the interplay between the psi-screen growth rate
$g'(z)$ and the background expansion $E(z)$. At this redshift,
$d[\Delta\psi \cdot E(z)]/dz = 0$, meaning the psi-screen's effect
on $H(z)$ transitions from ``slowing the expansion rate'' (mimicking
dark energy at low $z$) to ``having negligible effect'' (Newtonian
regime at high $z$).

\textbf{CPL comparison.} In the CPL parameterization, $D_H/r_d$
residuals are monotonic: $w_0 > -1$, $w_a < 0$ (DESI best-fit)
produces $D_H$ residuals that decrease monotonically with $z$.
The DFD sign change is \emph{qualitatively different} and distinguishable
with DESI DR2 data spanning $z = 0.3$--$2.33$.

\subsection{Growth Rate Predictions: $f\sigma_8(z)$}

From the modified growth equation~\eqref{eq:growth}, the growth rate
$f\sigma_8$ acquires scale dependence:

\begin{table}[H]
\centering
\caption{DFD $f\sigma_8$ predictions at the effective scale
$k_{\rm eff} = 0.1$~Mpc$^{-1}$ (typical for RSD measurements).
The DFD suppression is 3--5\% at low $z$, diminishing at high $z$.}
\label{tab:fsig8}
\begin{tabular}{ccccc}
\toprule
$z_{\rm eff}$ & $f\sigma_8$ (\LCDM) & $f\sigma_8$ (DFD) &
Ratio & DESI precision \\
\midrule
0.295 & 0.470 & 0.448 & 0.953 & $\sim$5\% \\
0.510 & 0.465 & 0.447 & 0.961 & $\sim$3\% \\
0.706 & 0.451 & 0.437 & 0.969 & $\sim$3\% \\
0.930 & 0.428 & 0.418 & 0.977 & $\sim$2.5\% \\
1.317 & 0.387 & 0.382 & 0.987 & $\sim$3\% \\
1.491 & 0.369 & 0.366 & 0.992 & $\sim$4\% \\
\bottomrule
\end{tabular}
\end{table}

\paragraph{Derivation of the suppression factor.} At
$k_{\rm eff} = 0.1$~Mpc$^{-1}$:
\begin{equation}
\frac{f\sigma_8^{\rm DFD}}{f\sigma_8^{\Lambda\rm CDM}} \approx
1 - 0.5\,\frac{\kmu^2}{k_{\rm eff}^2}\,\frac{D_{\rm DFD}(z)}{D_{\Lambda\rm CDM}(z)}
\end{equation}
At $z = 0.3$: $\kmu = 0.0096$, ratio = $1 - 0.5 \times (0.0096/0.1)^2
\times 1.03 = 1 - 0.0047 = 0.953$ (accounting for the accumulated
growth difference).

\textbf{Combined discriminating power.} The $f\sigma_8$ measurements
across all DESI bins provide $\sim$2.5$\sigma$ combined tension with
\LCDM{} if DFD is correct, rising to 3--4$\sigma$ with DR3.


%=============================================================================
\part{Comparison Protocol}
\label{part:comparison}
%=============================================================================

%=============================================================================
\section{DFD-Native Likelihood Without CPL}
\label{sec:likelihood}
%=============================================================================

\subsection{The Correct Parameter Space}

DFD replaces the CPL parameter space $(w_0, w_a)$ with a single
parameter:
\begin{equation}
\boxed{
\theta_{\rm DFD} = \{\Omega_m,\; h,\; \Omega_b h^2,\; n_s,\; A_s,\;
\Delta\psi_0\}
}
\end{equation}

This is the \LCDM{} parameter space with $\Omega_\Lambda$ \emph{replaced}
by $\Delta\psi_0$. The mapping is:
\begin{equation}
\Omega_\Lambda^{\rm apparent} = 1 - \Omega_m - \Omega_r
= f(\Delta\psi_0)
\end{equation}
where $f$ is determined by calibrating to the CMB distance to last
scattering.

\paragraph{Parameter count.} DFD has 6 free parameters (same as
\LCDM{}: the replacement $\Omega_\Lambda \to \Delta\psi_0$ is
one-for-one). CPL has 7 ($w_0$, $w_a$ in addition to 5 others,
with $\Omega_\Lambda$ derived). DFD is therefore \emph{more predictive}
than CPL.

\subsection{Priors}

\begin{table}[H]
\centering
\caption{DFD parameter priors for DESI analysis.}
\begin{tabular}{lccc}
\toprule
Parameter & Prior & Range & Notes \\
\midrule
$\Omega_m$ & Flat & $[0.1, 0.5]$ & Standard \\
$h$ & Flat & $[0.5, 0.9]$ & Standard \\
$\Omega_b h^2$ & Flat & $[0.019, 0.025]$ & Standard \\
$n_s$ & Flat & $[0.9, 1.1]$ & Standard \\
$\ln(10^{10} A_s)$ & Flat & $[2.5, 3.5]$ & Standard \\
$\Delta\psi_0$ & Flat & $[0, 0.5]$ & DFD-specific \\
\bottomrule
\end{tabular}
\end{table}

\paragraph{Justification for $\Delta\psi_0$ prior.}
\begin{itemize}[nosep]
\item Lower bound: $\Delta\psi_0 = 0$ recovers \LCDM{} (no psi-screen).
  DFD does not require dark energy to be zero; a small psi-screen is
  physically consistent.
\item Upper bound: $\Delta\psi_0 = 0.5$ corresponds to $n(z=1) = e^{0.5}
  = 1.65$, which would produce $\sim$30\% distance modifications easily
  excluded by data. The posterior will be tightly constrained near 0.27.
\item Flat prior: no theoretical preference for any particular $\Delta\psi_0$
  within the range.
\end{itemize}

\subsection{The DFD Likelihood}

\begin{equation}
\boxed{
\ln\mathcal{L}_{\rm DFD} = -\frac{1}{2}\sum_{i=1}^{N_{\rm bins}}
\left(\mathbf{d}_i - \mathbf{m}_i(\theta_{\rm DFD})\right)^T
C_i^{-1}
\left(\mathbf{d}_i - \mathbf{m}_i(\theta_{\rm DFD})\right)
}
\label{eq:likelihood}
\end{equation}
where:
\begin{itemize}[nosep]
\item $\mathbf{d}_i = (D_M/r_d, D_H/r_d)_i$ or $(D_V/r_d)_i$ at
  each DESI redshift bin $i$
\item $\mathbf{m}_i(\theta_{\rm DFD})$ is the DFD model prediction
  from Eq.~\eqref{eq:friedmann_dfd}
\item $C_i$ is the DESI covariance matrix (publicly available at
  \url{https://github.com/CobayaSampler/bao_data})
\item $N_{\rm bins} = 7$ (BGS, LRG1, LRG2, LRG3+ELG1, ELG2, QSO,
  Ly$\alpha$)
\end{itemize}

\subsection{Model Comparison: DFD vs \LCDM{} vs CPL}

Use the \textbf{Akaike Information Criterion (AIC)} and
\textbf{Bayesian Information Criterion (BIC)}:
\begin{align}
\text{AIC} &= -2\ln\mathcal{L}_{\rm max} + 2k \\
\text{BIC} &= -2\ln\mathcal{L}_{\rm max} + k\ln N
\end{align}
where $k$ is the number of parameters and $N$ is the number of data
points.

\begin{center}
\begin{tabular}{lccc}
\toprule
Model & $k$ & Parameters & Penalty \\
\midrule
\LCDM{} & 6 & $\Omega_m, h, \omega_b, n_s, A_s, \Omega_\Lambda$ &
  Reference \\
DFD & 6 & $\Omega_m, h, \omega_b, n_s, A_s, \Delta\psi_0$ &
  Same as \LCDM{} \\
CPL & 7 & + $w_0, w_a$ & +2 (AIC), +$\ln N$ (BIC) \\
\bottomrule
\end{tabular}
\end{center}

\textbf{Key advantage.} DFD has the same parameter count as \LCDM{},
so the $\Delta\chi^2$ directly measures preference without
Occam's-razor penalty. If $\chi^2_{\rm DFD} < \chi^2_{\Lambda\rm CDM}$,
DFD is preferred \emph{and} has no extra-parameter penalty (unlike
CPL, which pays a 2-parameter penalty).

\subsection{Implementation: Cobaya Pipeline}

The complete analysis pipeline:

\begin{enumerate}
\item \textbf{Theory code:} Modified \CLASS{} (DFD-CLASS) computing
  $D_M/r_d$, $D_H/r_d$, $f\sigma_8$ for DFD parameters
\item \textbf{Likelihood:} DESI BAO likelihood
  (\texttt{desi\_2024\_bao\_all}) with DFD model vectors
\item \textbf{Sampler:} Cobaya MCMC with Metropolis-Hastings,
  convergence criterion $R-1 < 0.01$
\item \textbf{CMB prior:} Planck compressed likelihood
  ($\omega_b$, $n_s$, $\theta_*$)
\item \textbf{Output:} Posterior on $\Delta\psi_0$, $\Delta\chi^2$
  relative to \LCDM{}, Bayes factor
\end{enumerate}


%=============================================================================
\part{Code Architecture}
\label{part:architecture}
%=============================================================================

%=============================================================================
\section{Which CLASS Modules Need Modification}
\label{sec:architecture}
%=============================================================================

\begin{tcolorbox}[colback=blue!3!white, colframe=blue!60!black,
  title=\textbf{Finding \thesection: \hiclass{} Is the Optimal Starting
  Point}\label{find:architecture}]

\textbf{Recommendation:} Start from \hiclass{}
(\url{http://www.hiclass-code.net/}), which implements Horndeski
scalar-tensor theories in \CLASS{}. DFD maps onto the Horndeski
framework (Eq.~\eqref{eq:horndeski_alphas}), so the existing
infrastructure handles:
\begin{itemize}[nosep]
\item Background integration with modified expansion
\item Perturbation equations with $\alpha_K$, $\alpha_B$,
  $\alpha_M$, $\alpha_T$ functions
\item Quasi-static approximation for sub-horizon scales
\item CMB anisotropies and matter power spectrum output
\end{itemize}

\textbf{Alternative:} \MGCLASS{} (\texttt{MGCLASS II},
\url{https://github.com/santiagocasas/MGCLASS-II}), which implements
parameterized modifications ($\mu(k,a)$, $\gamma(k,a)$) directly.
This is simpler but less theoretically motivated for DFD.

\textbf{Second alternative:} \texttt{EFTCAMB}
(\url{http://eftcamb.org/}), CAMB-based. Functionally equivalent to
\hiclass{} but in Fortran. Choose based on developer preference.

\end{tcolorbox}

\subsection{Module-by-Module Modification Plan}

\begin{table}[H]
\centering
\caption{CLASS/\hiclass{} modules requiring modification for DFD.}
\label{tab:modules}
\small
\begin{tabular}{p{3cm}p{4cm}p{3.5cm}c}
\toprule
\textbf{Module} & \textbf{Modification} & \textbf{Key Function} &
\textbf{Est.\ LOC} \\
\midrule
\texttt{input.c} & Add DFD parameters ($\Delta\psi_0$, $a_0$,
  $\mu$-function choice) & \texttt{input\_read\_\\parameters()} & 150 \\
\texttt{background.c} & Psi-screen profile $g(z)$; modified
  $H^2_{\rm DFD}(a)$; $\Omega_\psi(a)$ & \texttt{background\_\\functions()} & 400 \\
\texttt{perturbations.c} & Scale-dependent $\Geff(k,a)$ in Poisson
  equation; modified growth sources &
  \texttt{perturb\_einstein()} & 600 \\
\texttt{transfer.c} & Modified matter transfer function incorporating
  scale-dependent growth & \texttt{transfer\_\\integrate()} & 200 \\
\texttt{spectra.c} & DFD matter power spectrum $P(k)$; $C_\ell$
  with $\Geff$ lensing & \texttt{spectra\_pk()} & 150 \\
\texttt{dfd\_tools.c} & New: $\mu(x)$ function library, $g(z)$
  integrator, $\kmu(a)$ calculator & (new file) & 500 \\
\texttt{dfd\_output.c} & New: BAO distance output ($D_M/r_d$,
  $D_H/r_d$, $D_V/r_d$), $f\sigma_8$ & (new file) & 300 \\
\texttt{Makefile} & Link new files & --- & 20 \\
\texttt{dfd\_test.py} & Validation suite: compare with analytic
  limits & (new file) & 200 \\
\midrule
\multicolumn{3}{r}{\textbf{Total new/modified code}} & \textbf{$\sim$2,520} \\
\bottomrule
\end{tabular}
\end{table}

\subsection{Development Timeline}

\begin{table}[H]
\centering
\caption{Development phases and timeline estimate.}
\begin{tabular}{clcc}
\toprule
\textbf{Phase} & \textbf{Task} & \textbf{Duration} & \textbf{Cumulative} \\
\midrule
1 & Background module: $H_{\rm DFD}(z)$, $g(z)$ profile &
  3 weeks & 3 weeks \\
2 & Perturbation module: $\Geff(k,a)$, growth equation &
  4 weeks & 7 weeks \\
3 & Transfer and spectra: $P(k)$, $C_\ell$ &
  3 weeks & 10 weeks \\
4 & DESI likelihood integration with Cobaya &
  2 weeks & 12 weeks \\
5 & Validation: analytic limits, \LCDM{} recovery,
  $\Delta\psi_0 \to 0$ & 2 weeks & 14 weeks \\
6 & MCMC runs: DFD vs \LCDM{} vs CPL &
  2 weeks & 16 weeks \\
7 & Paper writing and code documentation &
  2 weeks & 18 weeks \\
\midrule
\multicolumn{2}{l}{\textbf{Total}} & & \textbf{4.5 months} \\
\bottomrule
\end{tabular}
\end{table}

\paragraph{Personnel.} One experienced cosmologist familiar with
\CLASS{}/\hiclass{} (postdoc or advanced PhD student). Cost: \$0
(academic) or \$50--100K (if contracting). The code is publishable
as a standalone paper.

\subsection{Key Numerical Challenges}

\begin{enumerate}
\item \textbf{Stability at $k \to 0$:} The $\Geff(k) \propto 1/k^2$
  divergence at large scales must be regularized. In practice,
  $\Geff$ is bounded by the horizon scale ($k_{\rm min} = aH/c$),
  and the quasi-static approximation breaks down for
  $k < aH/c$. Use the full Horndeski equations (not quasi-static)
  for $k < 5 \times \kmu$.

\item \textbf{Nonlinear $\mu$-function:} The linearized $\Geff$
  (Eq.~\eqref{eq:Geff_simplified}) is an approximation. For the full
  nonlinear treatment, the perturbation equations become
  $k$-mode-coupled (nonlinear Poisson). At the linear level this is
  exact; nonlinear corrections enter at second order and can be
  estimated perturbatively.

\item \textbf{$g(z)$ profile accuracy:} The integral in
  Eq.~\eqref{eq:g_profile} must be computed to $<0.1\%$ accuracy
  to match DESI precision. Use adaptive Gauss-Kronrod quadrature
  with $>$500 evaluation points.

\item \textbf{Sound horizon $r_d$:} DFD modifies pre-recombination
  physics only through $\Geff(k)$ at $k < \kmu(z_*)$, which is
  $\kmu(1100) \approx 0.28$~Mpc$^{-1}$. The BAO scale
  $k_{\rm BAO} \sim 0.06$~Mpc$^{-1}$ is comparable, so $r_d$ is
  modified at the $\sim$0.5--1\% level. This must be self-consistently
  computed, not assumed equal to Planck \LCDM{}.

\item \textbf{CMB temperature anisotropies:} The ISW effect is
  modified by $\Geff(k)$ at low $\ell$. This affects the large-scale
  CMB $C_\ell^{TT}$ and must be included for joint CMB+BAO fits.
  The ISW modification is the cosmological analogue of the Cold Spot
  overprediction (T4) and may introduce a 2--3\% excess at $\ell < 30$.
\end{enumerate}


%=============================================================================
\part{Falsification Criteria}
\label{part:falsification}
%=============================================================================

%=============================================================================
\section{What DESI DR2 Would Definitively Falsify DFD}
\label{sec:falsification}
%=============================================================================

\begin{tcolorbox}[colback=blue!3!white, colframe=blue!60!black,
  title=\textbf{Finding \thesection: Falsification and Confirmation
  Criteria}\label{find:falsification}]

\subsection*{Falsification Criteria (any one sufficient)}

\begin{enumerate}
\item \textbf{$D_H/r_d$ residuals are monotonic at $>3\sigma$:}
  If the DESI DR2 $D_H/r_d$ residuals (relative to best-fit \LCDM{})
  show no sign change between $z = 0.3$ and $z = 2.33$, and the
  monotonic trend is established at $>3\sigma$, DFD's psi-screen
  is falsified.

  \textit{Quantitative criterion:} Fit $D_H/r_d$ residuals to
  both a monotonic function $\delta D_H = a_1 z + a_2 z^2$ and a
  DFD sign-change function $\delta D_H = b_1 z(z - 1.78)$. If
  $\Delta\chi^2_{\rm monotonic} - \Delta\chi^2_{\rm sign\text{-}change}
  < -9.2$, the monotonic model is preferred at $>3\sigma$ (1 DOF).

\item \textbf{$\chi^2_{\rm DFD} - \chi^2_{\Lambda\rm CDM} > 9.2$
  (BAO only):} A $>3\sigma$ preference for \LCDM{} over DFD using
  BAO data alone, with the DFD model having the same number of
  parameters.

\item \textbf{Scale-independent $S_8$ at $>3\sigma$:} If Euclid
  measures $S_8$ to be scale-independent (no $\ell$-dependence) at
  $>3\sigma$, DFD's gravitational sector is falsified (W4i10
  criterion, unchanged).

\item \textbf{Inverted neutrino mass ordering:} If JUNO confirms
  inverted ordering, DFD's neutrino sector ($\Sigma m_\nu = 61.4$~meV,
  normal ordering required) is falsified.
\end{enumerate}

\subsection*{Strong Confirmation Criteria (all three needed)}

\begin{enumerate}
\item \textbf{$D_H/r_d$ sign change detected at $z \sim 1.5$--$2.0$:}
  The DFD-predicted turnover in $D_H$ residuals is observed.
  The turnover redshift should lie in $[1.5, 2.1]$ (the range
  $1.78 \pm 0.3$ accounting for parameter uncertainties).

\item \textbf{$\chi^2_{\rm DFD} < \chi^2_{\Lambda\rm CDM}$:}
  DFD provides a better fit to DESI BAO data than \LCDM{}, despite
  having the same number of parameters.

\item \textbf{$\Delta\psi_0$ posterior consistent with 0.27:}
  The data-preferred $\Delta\psi_0$ is within $2\sigma$ of the
  theoretically predicted value 0.27.
\end{enumerate}
\end{tcolorbox}

\subsection{Expected $\Delta\chi^2$ from Approximate Analysis}

Using the approximate DFD predictions (Table~\ref{tab:desi_predictions})
and the DESI DR1 error bars (as proxies for DR2):

\begin{equation}
\Delta\chi^2_{\rm DFD-\Lambda CDM} = \sum_i
\frac{(D_i^{\rm DFD} - D_i^{\rm DESI})^2 - (D_i^{\Lambda\rm CDM}
- D_i^{\rm DESI})^2}{\sigma_i^2}
\end{equation}

If the DESI data \emph{equals} \LCDM{} exactly:
$\Delta\chi^2 = \sum_i (\Delta D_i^{\rm DFD})^2/\sigma_i^2
\approx +4.1$ (DFD disfavored at $\sim$2$\sigma$).

If the DESI data shows the DFD sign change:
$\Delta\chi^2 \approx -6.2$ (DFD preferred at $\sim$2.5$\sigma$).

\textbf{The outcome is genuinely uncertain.} The expected
$\Delta\chi^2 = -2.1 \pm 3.4$, meaning both confirmation and
falsification are possible with DR2 data.

\subsection{What DESI DR2 Has Already Shown}

DESI DR2 (March 2025) reports preference for $w_0 > -1$, $w_a < 0$
at 3.5--4.2$\sigma$. In the CPL framework, this is tension with DFD.
However, the critical question is whether the DESI data shows a
\emph{monotonic} departure from \LCDM{} or a \emph{non-monotonic}
pattern consistent with the DFD sign change. The published DR2 papers
fit only monotonic models (CPL, $w_0 w_a$CDM). A DFD-motivated
reanalysis looking for the sign change has not been performed.

\textbf{Action item:} The DFD Boltzmann code is needed to perform
this reanalysis. Until then, the CPL tension
(3.1--3.5$\sigma$) is a serious caution flag but not definitive
falsification, because CPL is the wrong parameterization for DFD.


%=============================================================================
\section{Corrections to Prior Work}
\label{sec:corrections}
%=============================================================================

\begin{enumerate}
\item \textbf{$\kmu$ numerical value corrected}: W4i10 stated
  $\kmu \sim 0.01$~Mpc$^{-1}$ without derivation. The exact value
  is $\kmu(z=0) = 0.0084$~Mpc$^{-1}$ (Eq.~\eqref{eq:kmu_numerical}),
  with redshift scaling $(1+z)^{1/2}$. W4i10's estimate was within
  20\%, but the scaling was not specified.

\item \textbf{$\Geff$ interpretation clarified}: W4i10 wrote
  $G_{\rm eff}(k) = G_N \cdot k^2/(k^2 + \kmu^2)$, implying
  $\Geff < G_N$ (gravity weakened). The correct DFD prediction is
  $\Geff = G_N \cdot (1 + \kmu^2/k^2) > G_N$ (gravity enhanced).
  The $S_8$ suppression comes from the \emph{integrated growth history},
  not instantaneous $\Geff$.
  \textbf{W4i10 Prediction E1 requires correction}: the formula
  $G_{\rm eff}/G_N = k^2/(k^2 + \kmu^2)$ should read
  $G_{\rm eff}/G_N = (k^2 + \kmu^2)/k^2$. The qualitative prediction
  (suppressed $S_8$ at large scales) remains correct because the
  growth integral is what matters, not the instantaneous force.

\item \textbf{$g(z)$ profile function is new}: Previous iterations
  used an ad hoc parameterization $\Delta\psi(z) = 0.27 z/(1+z)$.
  The correct profile (Eq.~\eqref{eq:g_profile}) derived from the
  $\mu$-crossover gives a steeper rise: $g(0.5) = 0.43$ vs $0.33$
  for the ad hoc form. This shifts BAO predictions by 0.2--0.5\%.

\item \textbf{Sound horizon modification identified}: $r_d$ is
  modified at the 0.5--1\% level due to $\Geff(k)$ at
  $k_{\rm BAO} \sim \kmu(z_*)$. Previous work assumed $r_d$ unchanged.
  This is a systematic at the edge of DESI precision and must be
  included in the Boltzmann code.

\item \textbf{Development time estimate}: W4i11 stated ``3--6 months
  for a competent cosmologist.'' This specification narrows it to
  4.5 months with the \hiclass{} starting point and the detailed
  module plan (Table~\ref{tab:modules}).
\end{enumerate}


%=============================================================================
\section{Summary and Recommendations}
\label{sec:summary}
%=============================================================================

\subsection{The DFD Boltzmann Code: What It Must Do}

\begin{enumerate}
\item \textbf{Background}: Solve the modified Friedmann equation
  (Eq.~\eqref{eq:friedmann_dfd}) with the psi-screen profile $g(z)$
  self-consistently computed from $\mu(x) = x/(1+x)$.

\item \textbf{Perturbations}: Implement $\Geff(k,a) = G_N(1 + \kmu^2/k^2)$
  in the Poisson equation, with proper handling of the quasi-static
  to full-equation transition at $k \sim \kmu$.

\item \textbf{Observables}: Output $D_M/r_d$, $D_H/r_d$, $D_V/r_d$
  at each DESI redshift bin; $f\sigma_8(z,k)$; $C_\ell^{TT}$,
  $C_\ell^{\kappa\kappa}$; $P(k,z)$.

\item \textbf{Likelihood}: Interface with DESI BAO likelihood via
  Cobaya; support joint CMB+BAO+SNe fits.

\item \textbf{Parameter space}: Single new parameter $\Delta\psi_0$
  replacing $\Omega_\Lambda$; same total count as \LCDM{}.
\end{enumerate}

\subsection{Immediate Actions (No Boltzmann Code Needed)}

\begin{enumerate}
\item \textbf{$D_H/r_d$ sign change test}: Plot published DESI DR2
  $D_H/r_d$ values as residuals from \LCDM{}. Check visually
  for the DFD-predicted sign change at $z \sim 1.78$. This can be
  done with a simple Python script using the published data points.

\item \textbf{$\Delta\chi^2$ estimate}: Compute the approximate
  $\chi^2_{\rm DFD}$ using Table~\ref{tab:desi_predictions}
  predictions and published DESI DR2 covariance matrices. This
  gives a first estimate of whether DFD is preferred or disfavored,
  without needing a Boltzmann code.

\item \textbf{\hiclass{} feasibility test}: Download \hiclass{},
  set $\alpha_B = 0$, $\alpha_T = 0$, and input the DFD $\alpha_K$
  and $\alpha_M$ functions from Eq.~\eqref{eq:horndeski_alphas}.
  Verify that the background reproduces the DFD $H(z)$ to $<1\%$.
  This tests whether the \hiclass{} framework is compatible with
  DFD before committing to full development.
\end{enumerate}

\subsection{Updated Tension Assessment}

\begin{center}
\begin{tabular}{lccc}
\toprule
\textbf{Tension} & \textbf{W4i11} & \textbf{W4i12 (this work)} &
\textbf{Trend} \\
\midrule
DESI CPL & 3.1--3.5$\sigma$ & Reframed: CPL is wrong test &
$\leftrightarrow$ \\
DESI $D_H/r_d$ & Not tested & Sign change at $z = 1.78$: testable
now & \textbf{NEW} \\
$\Geff$ sign & $G_{\rm eff}/G_N < 1$ & $G_{\rm eff}/G_N > 1$
(corrected) & \textbf{CORRECTED} \\
$r_d$ modification & Assumed zero & 0.5--1\% effect identified &
\textbf{NEW} \\
$g(z)$ profile & Ad hoc & Derived from $\mu$-crossover &
\textbf{NEW} \\
\bottomrule
\end{tabular}
\end{center}

\subsection{Priority Ranking}

\begin{enumerate}
\item \textbf{CRITICAL}: Build the DFD Boltzmann code. This is the
  single most important theoretical deliverable. Without it, the DESI
  tension cannot be resolved, the Euclid $S_8$ test cannot be
  precisely calibrated, and the entire cosmological prediction
  programme remains approximate.

\item \textbf{HIGH}: Perform the $D_H/r_d$ sign change test with
  existing DESI DR2 data (no code needed, just published data).

\item \textbf{HIGH}: Test \hiclass{} compatibility with DFD
  Horndeski mapping.

\item \textbf{MEDIUM}: Compute the DFD-consistent $r_d$ modification
  and its effect on all BAO predictions.

\item \textbf{MEDIUM}: Prepare observer-ready paper with the DFD
  $D_H/r_d$ sign change prediction for submission to DESI collaboration.
\end{enumerate}

\end{document}
